\section{Direct Action y Reward Program: comparación a nivel de sistema}

Las dos partes de la Lecture no se contradicen entre sí. RT-2 incorpora la action generation al modelo grande y obtiene representaciones compartidas e inference rápida; Language to Rewards conserva el optimizer y obtiene constraint handling y modificación interactiva de objetivos. La elección de arquitectura real depende de los datos, la complejidad de las acciones, la frecuencia de control, la disponibilidad de modelos y los requisitos de seguridad.

\subsection{Diapositiva de resumen final: modelo, datos e interfaz}

La diapositiva de Final summary resume la primera parte como model consolidation, joint scaling y positive transfer, y la segunda parte como new interfaces to LLMs. La Key takeaway vuelve a enfatizar: rethinking scaling puede hacer que la robotics se parezca más al foundation-model training; rethinking interface permite que el LLM se encargue de la semantic translation en lugar de asumir todos los detalles físicos de bajo nivel.

\lecturefigure{slide-54-final-summary.jpg}{Final summary: Language to Rewards y general pattern machines}{Diapositiva V054 recuperada del video oficial; 01:09:39}

\lecturefigure{slide-55-key-takeaway.jpg}{Key takeaway: escalar el modelo de forma conjunta y también rediseñar la interfaz}{Diapositiva V055 recuperada del video oficial; 01:10:39}

\subsection{Tradeoff table de las dos rutas}

La siguiente tabla no busca elegir un ganador permanente, sino dejar clara la component responsibility. Los sistemas también pueden ser híbridos: primero se usa Language to Rewards en el simulator para generar diverse trajectories y luego se hace distill hacia una VLA; o bien la VLA propone el skill y el MPC se encarga del último tramo de contacto de alta precisión.

\begin{longtable}{@{}L{0.20\textwidth}L{0.32\textwidth}L{0.32\textwidth}@{}}
\toprule
Dimensión & Direct action / VLA & Reward program + controller \\
\midrule
Salida principal & action tokens o parámetros de acción continuos & state/action reward code u objective \\
Origen semántico & web + robot co-fine-tuning & pretrained LLM + robot schema / prompt \\
Ejecución física & la policy aprende implícitamente la dynamics y el feedback & el optimizer usa explícitamente la dynamics y las constraint \\
Costo de inferencia & una o pocas pasadas de autoregressive decoding & cada paso requiere rollout / optimization \\
Requisitos de datos & gran cantidad de robot trajectories; admite distillation & puede escribir la reward sin ejemplos previos, pero requiere simulator/model y validation \\
Ventajas & de extremo a extremo, rápida, puede compartir representaciones con el VLM & acciones diversas, restricciones explícitas, objetivos modificables de forma interactiva \\
Riesgos principales & OOD action, quantization, behavior support, fallas no interpretables & reward hacking, code error, model bias, fallas de la optimización en línea \\
Escenarios adecuados & skill repetitivo de alta frecuencia, demonstrations suficientes, embodiment fija & tareas semánticas nuevas, constraint complejas, dynamics model disponible \\
\bottomrule
\end{longtable}

\begin{importantbox}{Esquema híbrido presentado en clase}
Usar Language to Rewards para generar simulator trajectories a gran escala y semánticamente diversas, y luego hacer distill de esos datos hacia un Transformer/VLA, permite aprovechar a la vez la capacidad de exploración del optimizer y la inference rápida de la policy. La dirección inversa también es viable: la VLA se encarga de la skill proposal de alto nivel y el MPC ejecuta localmente cumpliendo los límites de colisión, equilibrio o fuerza.
\end{importantbox}

\teachervoice{En la sesión de Q\&A, el ponente compara explícitamente las dos rutas: la direct action generation es muy coherente con el autoregressive modeling, pero puede tener dificultades para producir dexterous motion fuera de la distribución de entrenamiento; la reward generation aprovecha la optimization de la robotics tradicional, permite añadir safety constraint y produce acciones más diversas. Esta respuesta se acerca más al diseño de sistemas reales que la disyuntiva «de extremo a extremo o modular».}

\subsection{Los datos siguen siendo el mayor bottleneck}

Aunque la escala de los modelos crezca, el campo de los robots sigue careciendo de datos de lazo cerrado a escala de internet. En el lenguaje, la positive transfer ya es tan común que ni siquiera se destaca por separado, mientras que la robotics todavía busca señales tempranas. El cuello de botella de los datos no es solo de cantidad: la coverage de fallas, recuperaciones, estados peligrosos, tareas de largo plazo y distintas embodiment suele ser más escasa que la de demonstrations exitosas.

\begin{knowledgebox}{Seis preguntas para auditar los datos}
¿Cuántas trajectories independientes se recolectaron? ¿Se conservan las fallas? ¿Son diversos los entornos y los objetos? ¿La action distribution cubre la recovery? ¿Los distintos operator introducen style bias? ¿La evaluation se hace en environment y object realmente held-out? Si estas preguntas no tienen respuesta, reportar solo el «total de horas» no basta para juzgar la dataset quality.
\end{knowledgebox}

\subsection{La Physical safety no puede limitarse a heredar la LLM alignment}

Las últimas preguntas de los estudiantes se centraron en la seguridad. Los errores de un robot actúan directamente sobre las personas y el entorno, por lo que no se puede depender solo de que el LLM se niegue a responder. El sistema necesita hardware torque/force limit, emergency stop, workspace boundary, collision monitor, software interlock, safe controller, simulation test y human supervision. El ponente mencionó además que el equipo explora la constitutional safety, pero no dio detalles públicos; estas notas solo registran esa dirección futura y no presentan un sistema no publicado como una solución validada.

\begin{warningbox}{«El modelo no dice cosas peligrosas» no equivale a que el robot sea seguro}
Los objetivos de seguridad incluyen al menos: si la instrucción en sí es legítima, si el modelo entiende a qué objeto se refiere, si el plan entra en zonas peligrosas, si las acciones de bajo nivel exceden los límites de fuerza o velocidad, si el sistema hace fail safe cuando fallan los sensores y si el usuario puede detenerlo. La Language alignment cubre solo una parte; las safety layer de hardware y de controller deben existir de forma independiente.
\end{warningbox}

\teachervoice{El ponente dijo que el equipo se toma la safety «muy en serio», porque un sistema físico puede dañar a las personas o al entorno. Puso como ejemplo de límite de fuerza en hardware un robot finger desprendible, y usó la instrucción ambigua «meter al gato en el lavavajillas» para mostrar que un semantic misunderstanding puede convertirse en una catastrophic physical failure, y no en una simple hallucination de texto.}

\subsection{Resumen de la sección}

Las rutas VLA y reward-program asignan más responsabilidad a la learned policy o al optimizer, respectivamente. La primera destaca en representaciones compartidas e inference rápida; la segunda, en constraint explícitas y en la composición de tareas nuevas; los sistemas híbridos suelen ser más realistas. Sea cual sea la ruta, ni la robot data coverage ni las safety layers independientes pueden sustituirse con scale.
